%% kapa: exam=1 round=36 year=2025 period=2 form=A subject=appraisal_law subject_ko=감정평가관계법규 seq=28 number=28
%% source: Q-Net 감정평가사 기출문제 (cst00309) artlSeq=5231525 file=raw/1008/5231525/2025년 제36회 감정평가사 1차시험 2교시 문제지 원본.pdf
%% conversion: kordoc 4.24.51 -> tools/gen1.py (auto); review=proofread-sonnet; flags=none
\begin{question}{28}
공간정보의 구축 및 관리 등에 관한 법령상 지상 경계의 결정기준으로 옳지 않은 것은? (단, 다른 조건은 고려하지 않음)
\begin{choices}
  \item 연접되는 토지 간에 높낮이 차이가 없는 경우: 그 구조물 등의 중앙
  \item 연접되는 토지 간에 높낮이 차이가 있는 경우: 그 구조물 등의 상단부
  \item 도로ㆍ구거 등의 토지에 절토(땅깎기)된 부분이 있는 경우: 그 경사면의 상단부
  \item 토지가 해면 또는 수면에 접하는 경우: 최대만조위 또는 최대만수위가 되는 선
  \item 공유수면매립지의 토지 중 제방 등을 토지에 편입하여 등록하는 경우: 바깥쪽 어깨부분
\end{choices}
\end{question}
\answer{2}
